\subsection{Choosing a calibration scale}

The full-pool working point is supported by the held-out results in
\cref{tab:7-budget-full}; sampled alternatives are evaluated in
\cref{fig:7-budget-risk}. The independent-slot model in \cref{sec:6-2}
provides a planning scale: at $k=10$, $4.74/(Sk)\approx10^{-3}$ suggests
$S$ of order 500, while 128 has coarser resolution. This strong model
concerns reference disagreement, not GT loss, and does not make the
empirical recommendation a prespecified quality guarantee.

The path and common-order conditions for FULL are addressed in
\cref{sec:6-3}.


\subsection{Risk and selected-arm performance}

Over $m$ calibration draws, let draw $r$ select $\widehat\alpha_r$.
Define held-out loss and threshold-crossing frequency by
\begin{equation}
 L_r=\widehat r^*-\widehat r(\widehat\alpha_r),\qquad
 \widehat p_\delta=\frac{1}{m}\sum_{r=1}^{m}\mathbf1\{L_r>\delta\}.
 \label{eq:6-risk}
\end{equation}
Mean loss, crossing frequency and tail size answer different questions.
\Cref{tab:7-budget-quality,fig:7-budget-risk} report the archived summaries;
the largest supplementary per-draw losses are reported in \cref{sec:7-budget-trace}.

Reported selected-arm rate is
$\overline Q=m^{-1}\sum_r\widehat Q(\widehat\alpha_r)$, a lookup from timed
fixed arms. It is distinct from the median budget's rate and from the
throughput of a calibrated query stream. Quality and speed must be read
together (\cref{tab:7-budget-cost}).

\subsection{Cost of reuse}

For $B_{\rm reuse}$ equal-sized subsequent batches, model
\begin{equation}
 T_{\rm policy}(B_{\rm reuse})=T_{\rm cal}+B_{\rm reuse}T_{\rm chosen},\qquad
 T_{\rm fixed}(B_{\rm reuse})=B_{\rm reuse}T_f.
 \label{eq:6-total-cost}
\end{equation}
When $T_{\rm chosen}<T_f$, the break-even count is
\begin{equation}
 B_{\rm reuse}^*=\frac{T_{\rm cal}}{T_f-T_{\rm chosen}}.
 \label{eq:6-payback}
\end{equation}
Positive saving requires $T_{\rm chosen}<T_f$ and
$B_{\rm reuse}>B_{\rm reuse}^*$. Batch time is $T=1000n/Q$ milliseconds.
When choices vary, $1/\mathbb E[Q]\leq\mathbb E[1/Q]$, so inverse mean
QPS understates mean time. Timing boundaries are in
\cref{sec:a-budget-contract}; \cref{sec:7-stream-cost} tests direct
calibration-plus-reuse cost on four cells with quality alongside time.
Campaign estimates are in \cref{tab:7-budget-cost}.

Recalibration every $H$ batches adds $T_{\rm cal}/H$ per batch.
Drift detection and automatic recalibration remain future work.
